%% =====================================================================
\section{Discussion}\label{sec:discussion}
%% =====================================================================
\subsection{When Does the Confidence Pay Off?}
The per-frame confidence helps when updates are scarce, either because energy is scarce or because the classifier is weak. In both cases the UAV must choose which few disagreements to report, and the confidence ranks them. When energy is ample and the classifier is accurate, the UAV can report every disagreement, a wrong label is corrected within a slot or two, and a simple rule does as well as the belief. Three lessons for the design of IoT devices with on-board inference follow.

First, a device should never send a label that agrees with the receiver (Proposition~\ref{prop:confirm}). This rule alone separates the AoII-threshold policy and CA-DPP from the pre-identification and age-threshold policies, which differ from them by $52\%$ to $75\%$ at moderate budgets.

Second, the classifier should be calibrated on board. Isotonic calibration is a lookup table at run time, and it lets the belief use the confidence as a probability. The prior correction~\eqref{eq:lik} matters for unbalanced classes, which are the rule in monitoring: floods, fires and intruders are rare.

Third, the value of raw uploads depends on the link more than on the classifier. Uploads paid off at budgets of $2$\,mJ/slot and above, near the GS, and with small frames. At tight budgets the UAV almost never uploads, and the label-only version performs as well as the full policy.

\subsection{Correlated Errors}
The main weakness of CA-DPP at large budgets is Assumption~\ref{as:inf}(ii). The on-board errors in our traces come in runs (Fig.~\ref{fig:inference}(c)): an error is followed by another error with probability $0.56$, against a marginal error rate of $0.08$. The tempered update~\eqref{eq:post} discounts every frame equally, which helps on average but does not distinguish a new piece of evidence from a repetition of the last one. A belief that models the error process explicitly, for example with a hidden Markov model of the classifier state, could do so. It would enlarge the belief by a factor of two and change Step 1, but Lemma~\ref{lem:onestep} would still apply, since it depends only on $\mu_t$ and $S_t$. We see this as the most promising way to improve the policy at large budgets.

\subsection{Other IoT Settings}
Nothing in the model is specific to UAVs or floods. The link model~\eqref{eq:succ} and the energy model~\eqref{eq:energy} apply to any device whose channel varies predictably, such as a satellite in low Earth orbit that sees a ground station only during passes, or a ground sensor that wakes up on a schedule. For a satellite, long outages are predictable, and a version that spends more just before an outage, through the renewal formulation of Remark~\ref{rem:opt}, would anticipate them. A non-binary scene, such as a land-cover class, requires a belief over the AoII and the current class, but the structure of the gains carries over: an update removes AoII when the GS is wrong and the update right, and creates AoII in the opposite case.

\subsection{Limitations}
Five limitations remain. First, the policy minimizes the AoII of the next slot only, and we have no bound on its gap to the optimal policy of (P). Second, the threshold structure of Proposition~\ref{prop:threshold} is proved for the per-slot rule~\eqref{eq:rule} with a constant ground error. Whether the optimal policy of (P) has the same structure is open. Third, the Sentinel-2 imagery is a proxy for UAV frames. UAV frames have higher resolution, a different viewing geometry and motion blur, which would change both classifiers. Repeating the study with UAV datasets such as FLAME~\cite{shamsoshoara2021flame} and AIDER~\cite{kyrkou2020aider} is the next step. Fourth, the energy model is analytic. Measuring the energy of inference and transmission on a microcontroller would test whether the label really costs a tenth of an upload in practice. Fifth, the photo sequence and the UAV position are generated separately, so the imagery is not tied to the location on the circle. A spatial model of the area, with correlated neighboring patches, revisits and a route chosen by the UAV, is left for future work.

%% =====================================================================
\section{Conclusion}\label{sec:concl}
%% =====================================================================
We studied status updates from an IoT device whose sensor is its own classifier, with a UAV mapping a flood as the running example. The device tracks a belief over the AoII of the ground station's status, updated with the calibrated confidence of each frame. The belief gives closed-form one-step values of the three actions, which show that agreeing labels are worthless and that the per-slot decision has confidence thresholds ordered silent, raw, label. The online policy CA-DPP built on them meets the energy budget on every sample path and keeps a belief of about 1.2\,kB. On flood-mapping traces under chip-level cross-validation, it outperforms adaptations of three recent designs at every budget up to $5$\,mJ per slot. Against an AoII-threshold policy it helps when energy is scarce or the on-board classifier is weak, because then the confidence decides which disagreements are worth reporting. With an accurate classifier and ample energy, reporting every disagreement works as well. A belief that models the correlation of on-board errors, a performance bound through the renewal framework, experiments with UAV imagery and measured energy, and a spatial model of the area are the next steps.
